\section{The Propers: Themes and Movement}
\zlabel{triptych:concise:themes:start}
God's feast gives a life of communion: its guests must learn charity, and its abundance sustains them while they still know need. The long Gospel examines a guest already admitted. Paul thanks people whose help still matters despite his contentment. Between these questions, pardon, guidance, illumination and sacramental nourishment turn invitation into a life received from God and shared with others.

\subsection{Mercy opens the way to a changed life}
The Entrance's Psalm 130 begins with the sinner's inability to stand before a strict account of wrongdoing. Forgiveness makes approach possible. Augustine locates that hope in Christ's redemption and understands mercy as a companion on the continuing way. The Collect's concern with grace before and after good action gives this approach its direction. The faithful ask for help because there is good to be done, not because action has become unnecessary.

The garment, proclaimed only in the long Gospel, sharpens this relation. Gregory identifies charity as the life fitting to the Son's wedding. Augustine's account of purification explains how real change can depend on God's help: a believer acts, yet not alone. Mercy neither rewards a previously flawless guest nor excuses a settled refusal to love. The first reading of the formulary therefore asks what the gift of entry is becoming in those who have received it.

\subsection{The offered feast can be refused}
Isaiah promises a feast for all peoples, the destruction of death and the wiping away of tears. His mountain follows the vision of the Lord reigning in Zion; the surrounding chapter also speaks of refuge and judgment. In Matthew the king prepares and repeatedly sends before any guest is examined. Some neglect the invitation for farm and business; others attack the messengers. Those found on the roads, bad and good alike, then fill the room.

Both Gospel lengths proclaim that gathering. Gregory understands Matthew's wedding as the present Church because, in the long form, someone can still be expelled; the completed heavenly feast does not lose its guests. The assembly therefore cannot identify being inside with having reached the end. Nor can it demand a gathering emptied of everyone who needs conversion. Gregory makes endurance of another's faults a test of the charity the guest himself requires.

The acclamation's adaptation of Ephesians asks for the heart's illumination concerning the hope of the call. Chrysostom places that hope beside God's power in the risen Christ; Thomas distinguishes present hope from future inheritance. The invitation offers more than another possession. A heart absorbed in possessions or status needs sight to recognize what it is being offered and how to answer.

\clearpage
\subsection{Provision includes sharing another's distress}
The Responsorial Psalm's table stands amid enemies; the Shepherd accompanies the valley. God has not become absent because danger remains. Augustine reads the Church's growth through baptismal refreshment, right paths, discipline and nourishment. Thomas distinguishes the tables of instruction, sacrament and heaven. Their readings give present food a destination without claiming that the journey's troubles have already ended.

Paul makes the unfinished condition unmistakable: he knows hunger as well as plenty. Chrysostom explains that both require virtue, and that strength comes from Christ. Yet Paul immediately commends the Philippians' fellowship in his distress. Thomas draws the consequence: knowing how to bear want does not make help dispensable. The second reading of the formulary asks what God's abundance means for precisely such a person, whose trust and need are both real.

The Prayer over the Offerings relates finite offerings and prayer to heavenly glory. Paul's wider gift-context likewise distinguishes what is consumed from the fruit generosity bears. The offering cannot buy admission or replace love of a neighbor; its good exceeds the possession surrendered. A donor gives without owning the recipient, and a recipient receives without making the donor the source of salvation. The bodily act can become an act of communion.

\subsection{Nourishment tends toward a fullness still awaited}
The Communion alternatives give this movement two endings. The psalm option promises good to seekers of the Lord. Augustine tests that promise against a wealthy wrongdoer and a poor believer, then identifies the heart's true riches in justice, charity, faith, endurance and Christ the living bread. The Johannine option names likeness through future sight. Augustine describes present childhood as real while desire and purification enlarge the heart for a gift not yet fully manifested.

These emphases are complementary, but the antiphons are alternatives: the congregation does not hear them as one combined text. The Prayer after Communion joins sacramental nourishment with participation in divine life. Thomas's distinction between the Church's banquet, the soul's banquet and heaven gives this participation its breadth. The gift is present; the defeat of death promised by Isaiah remains the measure of its completion.

The two questions meet in practical charity. A garment without deeds would leave Paul's need untouched; a full table without charity would leave the Gospel's guest unchanged. Yet their difficulties differ. Gregory searches the person already inside for the life of love. Augustine and Chrysostom ask how to recognize God's good amid continued poverty and why help remains necessary. Both lead toward communion through grace, rather than through entitlement to either religious security or material success.
\zlabel{triptych:concise:themes:end}
\clearpage
